\section{A planted Gaussian comparison}
\label{rcut:sec:comparison}

We seek a bound on the quadratic form in
Equation~\eqref{rcut:eq:gradient-form} outside a set of configurations
of exponentially small Gibbs mass.  A bound of the form
$M_{J^x}(p)<\|p\|_2^2$ need not hold for every $x$ and every nonzero
$p$.  We therefore first fix a configuration and weight the disorder
by its unnormalized Gibbs weight.  Under this planted law, a scalar
variational estimate controls the form simultaneously on all unit
vectors.  The final step transfers the resulting exponential probability
bound to Gibbs mass under the original disorder law.

\subsection{The planted law and the variational bound}

Write $\widehat\P_n$ for the law under which the entries $J_{ij}$,
$i<j$, are independent Gaussians with mean $\beta^2/n$ and variance
$\beta^2/n$, with $J_{ji}=J_{ij}$ and $J_{ii}=0$.
Relative to the original disorder law $\P_n$, its density is
\begin{equation}
 \frac{d\widehat\P_n}{d\P_n}(J)
 =\exp\left\{\sum_{i<j}J_{ij}-\frac{\beta^2(n-1)}4\right\}.
 \label{rcut:eq:comp-planted-density}
\end{equation}
The density is the unnormalized Gibbs weight of the all-positive
configuration divided by its expectation.  Its normalization is therefore
deterministic; Equation~\eqref{rcut:eq:comp-planted-density} contains no
random partition function.

For $\beta>0$, let $\mathsf G_\beta=N(\beta^2,\beta^2)$ and write
$\langle f\rangle=\int f(h)\,d\mathsf G_\beta(h)$.
An asterisk denotes a second argument integrated against the same law.
In this Section the scalar functions are
\begin{equation}
 m(h)=\tanh h,\qquad v(h)=1-m(h)^2,\qquad w(h)=1-m(h),
 \label{rcut:eq:comp-field-functions}
\end{equation}
and we put
\begin{equation}
 D(h)=\beta^2\left\langle
             \frac{v(h_*)^2}{w(h)+w(h_*)}\right\rangle_*,
 \qquad
 K_0(h,h_*)=\frac{v(h)v(h_*)}{w(h)+w(h_*)}.
 \label{rcut:eq:comp-kernels}
\end{equation}
To control large fields uniformly, express the kernels in terms of
$w\in[0,2]$, using $v=w(2-w)$ also at the endpoints.  The resulting
kernels
\begin{equation}
 F(a,b)=\frac{b^2(2-b)^2}{a+b},\qquad
 K(a,b)=\frac{a(2-a)b(2-b)}{a+b}
 \label{rcut:eq:comp-compact-kernels}
\end{equation}
extend continuously to $[0,2]^2$ by taking value zero at $(0,0)$.
Indeed $F(a,b)\le b(2-b)^2\le32/27$, and
$K(a,b)\le4ab/(a+b)\le4\min(a,b)$.
Thus the denominators remain harmless after compactifying the field
through $w$, including the limit of large positive fields.

The variational parameters describe the part of a unit vector that is
constant at a given field and its remaining variation.  More precisely,
a finite partition into field bins gives a mean $b_a$ of the rescaled
coordinates $\sqrt n\,p_i$ and a root-mean-square deviation $r_a\ge0$
in each bin.  Equation~\eqref{rcut:eq:comp-bin-decomposition} makes this
decomposition precise.  With empirical bin proportions as weights, the
sum of $b_a^2+r_a^2$ is one.  The functional below uses the limiting
Gaussian law $\mathsf G_\beta$ instead; the rescaling in
Equation~\eqref{rcut:eq:comp-mass-rescaling} will preserve the unit norm
when we pass between the two measures.  Its final term will bound the
conditional Gaussian contribution from the within-bin variation.

For measurable real $b$ and nonnegative $r$ with
$\langle b^2+r^2\rangle=1$, define
\begin{align}
 \mathcal V_\beta(b,r)
 ={}&\langle D(b^2+r^2)\rangle
       +\langle bh\rangle\langle bv\rangle
       +\langle b\rangle\langle bhv\rangle
       -\beta^2\langle b\rangle\langle bv\rangle \notag\\
    &+2\beta^2\langle b\rangle\langle bmv\rangle
       +\beta^2\langle b(h)K_0(h,h_*)b(h_*)\rangle_{h,h_*}
       \notag\\
    &+\beta\left\langle r(h)
       \sqrt{\left\langle(v(h)+v(h_*))^2r(h_*)^2\right\rangle_*
        +\Var_{h_*}\bigl((v(h)+v(h_*))b(h_*)\bigr)}
       \right\rangle.
 \label{rcut:eq:variational}
\end{align}
This definition applies to the entire stated $L^2$ class.  Boundedness
of the kernels and of $m,v$, together with
$h,hv\in L^2(\mathsf G_\beta)$, makes the first two lines finite.
The last radical is bounded by
$\sqrt{4\langle b^2+r^2\rangle}=2$, so its pairing with $r$ is finite
as well.  In particular, $b$ and $r$ need not be bounded.

\begin{proposition}[Planted comparison]
 \label{rcut:prop:planted-comparison}
 Fix $0<\beta\le1/2$.  Suppose that, for some $\delta>0$,
 \[
  \sup_{\langle b^2+r^2\rangle=1,\ r\ge0}
       \mathcal V_\beta(b,r)\le1-\delta.
 \]
 There are $\kappa>0$ and $n_0<\infty$ such that for $n\ge n_0$,
 \begin{equation}
  \widehat\P_n\left(
    \sup_{\|p\|_2=1}M_J(p)>1-\frac\delta2\right)
       \le e^{-\kappa n}.
  \label{rcut:eq:comp-exponential-conclusion}
 \end{equation}
 The constants may depend on the fixed $\beta$ and $\delta$.
\end{proposition}

{\emergencystretch=3em
The proof separates the squared interactions from the centered Gaussian
matrix.  Lemma~\ref{rcut:lem:comp-square-replacement} replaces the squares
by their means without discarding the dependence between fields and
couplings.  After this replacement, a finite partition of the fields
reduces the matrix term to the conditional comparison in
Lemma~\ref{rcut:lem:comp-bin-gaussian}.  A net of the normalized bin
parameters then gives simultaneous control of all unit vectors.  We
choose every approximation accuracy before sending $n$ to infinity,
which retains the exponential failure rate required by the proposition.\par}

\subsection{Gaussian squares and the added diagonal}

The first task is to control weighted sums of interaction squares.
We begin with deterministic weights; the later conditioning argument
will explain when these estimates also apply to weights depending on
the fields.  Under $\widehat\P_n$,
$nJ_{ij}^2$ has the same law as
\[
 X_n=(\beta Z+\beta^2/\sqrt n)^2,\qquad Z\sim N(0,1).
\]
For $\mu=\beta^2/\sqrt n$ and $|t|<1/(2\beta^2)$, direct Gaussian
integration gives
\begin{align}
 \log\E e^{t(X_n-\E X_n)}
  ={}&-\tfrac12\log(1-2\beta^2t)-\beta^2t
      +\frac{2\beta^2\mu^2t^2}{1-2\beta^2t}
      \le C_\beta t^2,
 \label{rcut:eq:comp-square-mgf}
\end{align}
where the inequality holds for $|t|\le t_\beta$ with a fixed
$t_\beta>0$, uniformly in $n$.
Consequently, if $X_{n,j}$ are independent copies and $|a_j|\le A$,
then exponential Markov inequalities for both signs imply
\begin{equation}
 \P\left(\left|\frac1n\sum_{j=1}^{n-1}
       a_j(X_{n,j}-\E X_{n,j})\right|>u\right)
 \le 2\exp\{-c_{\beta,A}n\min(u^2,u)\}.
 \label{rcut:eq:comp-weighted-square-tail}
\end{equation}
To see the rate, independence bounds the logarithm of the moment
generating function of the unnormalized sum by $C_{\beta,A}nt^2$ on
a fixed interval of $t$.  In the exponential Markov bound choose $t$
to be a sufficiently small constant multiple of $\min(u,1)$, and repeat
with the opposite sign.  This argument remains valid after conditioning
on a sigma-field that determines the coefficients and is independent
of the copies.

Add independent diagonal entries $\zeta_i\sim N(\beta^2/n,2\beta^2/n)$
to the planted probability space, and set
\begin{equation}
 \widetilde J=J+\diag(\zeta_i)=W+\beta^2ee^{\mathsf T},\qquad
 e=\frac{\boldsymbol 1}{\sqrt n},\quad Q=I-ee^{\mathsf T},\quad
 H=\widetilde J\boldsymbol 1.
 \label{rcut:eq:comp-goe-coupling}
\end{equation}
Here $W$ is a centered GOE with off-diagonal variance $\beta^2/n$
and diagonal variance $2\beta^2/n$.  For every fixed $u>0$, Gaussian
tails, Equation~\eqref{rcut:eq:comp-square-mgf}, and
Lemma~\ref{rcut:lem:gaussian-facts} give constants $c,C>0$ such that,
outside an event of probability at most $Cn^2e^{-cn}$,
\begin{equation}
 \max_i|\zeta_i|\le u,\quad
 \max_{i<j}|J_{ij}|\le u,\quad
 \max_i\sum_jJ_{ij}^2\le\beta^2+u,\quad
 \|W\|\le2\beta+u.
 \label{rcut:eq:comp-basic-events}
\end{equation}
Each estimate has an exponential rate at fixed $u$.  A row sum has
mean $(n-1)(\beta^2/n+\beta^4/n^2)$, so
Equation~\eqref{rcut:eq:comp-weighted-square-tail} followed by a union
over rows gives the third bound.  Gaussian tails and finite unions give
the first two.  For the last, the GOE comparison in the proof of
Proposition~\ref{rcut:prop:good-disorder} gives
$\E\lambda_{\max}(W)\le2\beta$, with Gaussian Lipschitz constant
$\sqrt2\beta/\sqrt n$.  Apply Lemma~\ref{rcut:lem:gaussian-facts} to
both $W$ and $-W$.  These events also bound the planted norm $\|J\|$
by $2\beta+\beta^2+2u$.  Only a fixed bound on that norm is needed
in what follows.

The following decomposition is exact:
\begin{equation}
 \widetilde J
  =QWQ+\frac{He^{\mathsf T}+eH^{\mathsf T}}{\sqrt n}
       -ee^{\mathsf T}\frac{e^{\mathsf T}H}{\sqrt n}.
 \label{rcut:eq:comp-field-decomposition}
\end{equation}
This decomposition separates the field-dependent terms from a matrix
that is independent of $H$.  Indeed, the GOE covariance identity
\begin{equation}
 \E[(a^{\mathsf T}Wb)(c^{\mathsf T}Wd)]
  =\frac{\beta^2}{n}
     \bigl((a^{\mathsf T}c)(b^{\mathsf T}d)
                  +(a^{\mathsf T}d)(b^{\mathsf T}c)\bigr)
 \label{rcut:eq:comp-goe-covariance}
\end{equation}
makes the covariance of every entry of $QWQ$ with every entry of
$We$ equal to zero.  Joint Gaussianity then gives independence of
$QWQ$ and $H$.  The displayed decomposition itself follows by
expanding $W$ with $I=Q+ee^{\mathsf T}$.  The covariance identity
also gives the following representation of the field vector:
\begin{equation}
 H_i=\beta^2+\beta Z_i+\frac\beta{\sqrt n}Z_0,
 \label{rcut:eq:comp-field-law}
\end{equation}
where $Z_0,Z_1,\ldots,Z_n$ are independent standard Gaussians.
Indeed its covariance has diagonal $\beta^2(1+1/n)$ and off-diagonal
$\beta^2/n$.

\subsection{Replacing squared entries with their means}

\begin{lemma}[Weighted square replacement]
 \label{rcut:lem:comp-square-replacement}
 On the enlarged planted space, put $w_i=w(H_i)$ and
 $E_{ij}=J_{ij}^2-\beta^2/n$, including $E_{ii}=-\beta^2/n$.
 Here $\circ$ denotes entrywise product, and $F$ is defined in
 Equation~\eqref{rcut:eq:comp-compact-kernels}.
 For any fixed continuous function $G$ on $[0,2]^2$ and any $u>0$,
 there are $c>0$ and $n_0$ such that, for $n\ge n_0$, each of the
 following failures has probability at most $e^{-cn}$:
 \begin{align}
  &\left\|J\circ J-\beta^2ee^{\mathsf T}\right\|>u,
    \label{rcut:eq:comp-square-operator}\\
  &\sup_{\|p\|_2=1}
     \left|\sum_{i,j}E_{ij}G(w_i,w_j)p_ip_j\right|>u,
    \label{rcut:eq:comp-weighted-operator}\\
  &\max_i\left|\sum_{j\ne i}
       E_{ij}F(w_i,w_j)\right|>u.
    \label{rcut:eq:comp-weighted-diagonal}
 \end{align}
\end{lemma}

\begin{proof}
 We first establish the unweighted operator estimate, then insert the
 continuous kernel $G$ by approximation.  The diagonal estimate will
 require a separate argument that removes one row from the fields.
 For the operator estimate, truncate $X_n=nJ_{ij}^2$ at a fixed level $T$.
 For a sufficiently small fixed $t>0$,
 \[
  \sup_n\E e^{t(X_n-T)_+}\longrightarrow1
       \quad\text{as }T\longrightarrow\infty.
 \]
 For uniform dominated convergence, use
 $|\beta Z+\beta^2/\sqrt n|\le\beta|Z|+\beta^2$; the square of
 the right side has an exponential moment for a sufficiently small
 parameter.  We may therefore choose $T$ so that both the mean tail
 and the logarithm of the displayed moment are as small as required.
 Exponential Markov bounds the probability that a row average of
 $(nJ_{ij}^2-T)_+$ exceeds the prescribed accuracy by $e^{-cn}$.
 Taking the union over the $n$ rows retains a positive exponential
 rate.  Symmetry and entrywise nonnegativity bound the tail matrix's
 operator norm by its largest row sum, and the same estimate controls
 its expectation.

 For a fixed unit vector, the quadratic form of the centered truncated
 matrix is a sum of independent bounded variables with coefficients
 $2p_ip_j/n$, $i<j$.  To bound its moment generating function, recall
 that a centered variable whose range has length $a$ has logarithmic
 moment generating function at most $a^2s^2/8$.  The second derivative
 is a variance under a tilted law and is at most $a^2/4$; integrating
 twice gives this bound.  Since
 $\sum_{i<j}(2p_ip_j)^2\le2$, the present quadratic form has logarithmic
 moment generating function at most $CT^2s^2/n^2$, and hence a tail
 at fixed accuracy bounded by $2e^{-c n^2/T^2}$.
 To obtain an operator estimate, use a $1/4$-net of the unit sphere
 with at most $9^n$ points.  This cardinality follows by comparing
 volumes of disjoint balls of radius $1/8$.  For a symmetric matrix
 $A$, twice the maximum of $|p^{\mathsf T}Ap|$ on this net bounds
 $\|A\|$.  The union bound is therefore absorbed by the quadratic
 exponential rate.  It remains to restore the mean and diagonal:
 \[
  \E J_{ij}^2=\frac{\beta^2}{n}+\frac{\beta^4}{n^2}
       \quad(i\ne j),
 \]
 so the resulting mean correction and missing diagonal have norm
 $O_\beta(1/n)$.  Together with the truncation and net bounds, this
 proves Equation~\eqref{rcut:eq:comp-square-operator}.

 We now insert $G$.  Uniform continuity gives an approximation on a
 finite rectangular partition of $[0,2]^2$, of the form
 $\sum_{a,d}g_{ad}\1_{I_a}(w_i)\1_{I_d}(w_j)$.
 Multiplication of $p$ by the corresponding coordinate masks reduces
 each summand to Equation~\eqref{rcut:eq:comp-square-operator}.  The
 number of summands and their coefficients are fixed.  For the
 remaining kernel error $R_{ij}$, use
 \[
  \left|\sum_{i,j}E_{ij}R_{ij}p_ip_j\right|
   \le\max_{i,j}|R_{ij}|\,
          |p|^{\mathsf T}|E|\,|p|
   \le\max_{i,j}|R_{ij}|\max_i\sum_j|E_{ij}|.
 \]
 The row-sum bound in Equation~\eqref{rcut:eq:comp-basic-events}
 controls this error, since
 $\sum_j|E_{ij}|\le\sum_jJ_{ij}^2+\beta^2$.
 This proves Equation~\eqref{rcut:eq:comp-weighted-operator}.  On the
 operator and row-sum events, the argument is deterministic for every
 choice of masks and residual coefficients.  Thus these weights may
 depend on the whole interaction matrix.

 For the diagonal estimate we must recover independence before applying
 the weighted tail bound.  Fix a row $i$ and condition on
 \[
  \mathcal F_i=\sigma\bigl(J_{kl}:k,l\ne i;\ \zeta_1,\ldots,\zeta_n\bigr).
 \]
 This sigma-field leaves the entries of the off-diagonal row
 $(J_{ij})_{j\ne i}$ independent.  Remove that row from the other
 fields by defining
 \[
  H_j^{(-i)}=\zeta_j+\sum_{k\ne i,j}J_{jk},\qquad j\ne i.
 \]
 The fields so obtained are $\mathcal F_i$-measurable, with
 $H_j-H_j^{(-i)}=J_{ij}$.  Keep $a\in[0,2]$ deterministic for now
 and set $c_j(a)=F(a,w(H_j^{(-i)}))$.  Under the conditioning these
 weights are fixed and lie in $[0,32/27]$, so
 Equation~\eqref{rcut:eq:comp-weighted-square-tail} gives a bound
 $e^{-cn}$ for the event
 \[
  \left|\sum_{j\ne i}c_j(a)
      \left(J_{ij}^2-\frac{\beta^2}{n}\right)\right|>u,
 \]
 for all large $n$, uniformly in $\mathcal F_i$ and $a$; the mean
 correction is at most $(32/27)\beta^4/n$.
 Take a union first over a fixed finite mesh of $a$ values and then
 over the $n$ rows.  Uniform continuity of $F$ and the bound on
 $\sum_j|E_{ij}|$ extend the estimate to all $a\in[0,2]$.
 Next restore the original fields.  On the maximum-entry event,
 $|H_j-H_j^{(-i)}|$ is uniformly small.  The Lipschitz property of
 $w$ and uniform continuity of $F$ bound the resulting error by a
 small multiple of the same row sum.  Only after these simultaneous
 estimates do we set $a=w_i$.  Conditioning on $w_i$ at the outset
 would condition on a row sum and destroy the independence used above.
 The mesh, continuity tolerances and entry bound are chosen before $n$;
 their failures therefore retain positive exponential rates.  This
 completes Equation~\eqref{rcut:eq:comp-weighted-diagonal}.
\end{proof}

We can now simplify the quadratic form.  Before applying the replacement
lemma, change the fields in Equation~\eqref{rcut:eq:gradient-form} from
$J\boldsymbol1$ to $H$.  Each coordinate changes by $\zeta_i$.
For the squared-entry terms, uniform continuity of the compact kernels
in Equation~\eqref{rcut:eq:comp-compact-kernels} and the bounded row sums
make the error uniformly small on the unit sphere.  The unsquared term
has error at most
$\|J\|\max_i|v((J\boldsymbol1)_i)-v(H_i)|$; adding its diagonal
costs at most $\max_i|\zeta_i|$.  We then apply
Lemma~\ref{rcut:lem:comp-square-replacement} to the squared entries.

Put $A=\diag(v(H_i))$, $\overline H=n^{-1}\sum_iH_i$, and
\[
 D_n(t)=\frac{\beta^2}{n}\sum_j
       \frac{v(H_j)^2}{w(t)+w(H_j)}.
\]
For every fixed accuracy $u>0$, the preceding estimates show that,
outside probability $e^{-cn}$,
\begin{equation}
 \sup_{\|p\|_2=1}|M_J(p)-\widehat M(p)|\le u,
 \label{rcut:eq:comp-reduced-error}
\end{equation}
where
\begin{align}
 \widehat M(p)={}&p^{\mathsf T}QWQAp
 +\frac{p^{\mathsf T}H}{\sqrt n}(e^{\mathsf T}Ap)
 +(p^{\mathsf T}e)\frac{H^{\mathsf T}Ap}{\sqrt n}
 -\overline H(p^{\mathsf T}e)(e^{\mathsf T}Ap)\notag\\
 &+2\beta^2(p^{\mathsf T}e)
       \bigl(e^{\mathsf T}\diag(m(H_i)v(H_i))p\bigr)
 +\sum_iD_n(H_i)p_i^2
 +\frac{\beta^2}{n}\sum_{i,j}K_0(H_i,H_j)p_ip_j.
 \label{rcut:eq:comp-reduced-form}
\end{align}
Here the two weighted replacement estimates have different roles.
Equation~\eqref{rcut:eq:comp-weighted-diagonal} treats the diagonal part
of the last square in Equation~\eqref{rcut:eq:gradient-form};
Equation~\eqref{rcut:eq:comp-weighted-operator} treats its cross part
and the other squared-entry term.  The added terms with $i=j$ in
Equation~\eqref{rcut:eq:comp-reduced-form} contribute only
$O_\beta(1/n)$ on the unit sphere.  To check the coefficient of the
kernel term, expand the square over $i<j$: the cross part is exactly
the ordered off-diagonal sum with coefficient $\beta^2/n$.

\subsection{Finite field bins and normalized parameters}

The reduced form still depends on all $n$ field values.  We next
replace them by a fixed finite collection while retaining uniformity
over every unit vector.  Choose a finite partition
$\mathcal I=(I_1,\ldots,I_k)$ of the real line into intervals, with
the two extreme intervals unbounded, and choose a representative
$h_a\in I_a$ for each interval.  Write
\[
 h^{\mathcal I}(h)=h_a\quad(h\in I_a),\qquad
 m_a=m(h_a),\quad v_a=v(h_a),\quad w_a=w(h_a).
\]
For every prescribed accuracy, the partition and representatives can
be chosen so that $\|h-h^{\mathcal I}\|_{L^2(\mathsf G_\beta)}$ is
small, and so that $m,v$ and the kernels in
Equation~\eqref{rcut:eq:comp-compact-kernels} change uniformly by a
small amount on replacing each field by its representative.
Construct the partition by first choosing $R$ so that the Gaussian
square tails outside $[-R,R]$ are small.  In each tail take a
representative far enough out that $\tanh$ is uniformly close to its
endpoint; uniform continuity on the compactified $w$-square then
controls both kernels.  A fine mesh of $[-R,R]$ gives the remaining
approximations.  Choose every interval to have positive
$\mathsf G_\beta$-probability.

Let
\[
 \pi_a=\mathsf G_\beta(I_a),\qquad
 \widehat\pi_a=\frac1n|\{i:H_i\in I_a\}|.
\]
The following properties hold with exponentially small failure at
any fixed accuracy:
\begin{equation}
 \max_a|\widehat\pi_a-\pi_a|\ \text{is small},\qquad
 \frac1n\sum_i|H_i-h^{\mathcal I}(H_i)|^2\ \text{is small},
 \qquad |\overline H-\beta^2|\ \text{is small}.
 \label{rcut:eq:comp-empirical-events}
\end{equation}
We verify the exponential rates because the final comparison needs
more than convergence in probability.  In the representation
\eqref{rcut:eq:comp-field-law}, the common shift lies in any prescribed
small fixed interval except with probability $e^{-cn}$.  For the
independent fields $Y_i=\beta^2+\beta Z_i$, the bounded-variable
moment generating function estimate from the square-replacement proof
gives exponential concentration of each bin count.
A common shift of magnitude at most $a$ can move a field to a different
bin only when $Y_i$ is within $a$ of a boundary.  There are finitely
many boundaries.  Choose $a$ so that their neighborhoods have small
total probability, and concentrate their empirical count by the same
estimate.
For the empirical square error, the mesh, common shift and boundary
counts control the bounded middle interval.  The tail contribution
is at most a constant times the empirical average of
$(1+Y_i^2)\1_{\{|Y_i|>R-a\}}$.
These variables possess an exponential moment at a small fixed
parameter.  Increasing $R$ makes its logarithm arbitrarily small,
by the same domination argument used for Gaussian square tails.
Exponential Markov therefore makes the tail average small with
probability $1-e^{-cn}$.  The remaining assertion follows from
$\overline H-\beta^2\sim N(0,2\beta^2/n)$.
Thus every tolerance in Equation~\eqref{rcut:eq:comp-empirical-events}
has a fixed positive exponential rate.  These events also bound the
empirical second moment of $H$ and ensure, for large $n$, at least
two sites in every bin.

For a unit vector $p$, let $p_c$ be its orthogonal projection onto
vectors constant on every bin and write $p=p_c+u$.  On bin $a$ put
\begin{equation}
 (p_c)_i=\frac{b_a}{\sqrt n},\qquad
 \sum_{i:H_i\in I_a}u_i=0,\qquad
 \|u_a\|_2^2=\widehat\pi_a r_a^2,
 \quad r_a\ge0.
 \label{rcut:eq:comp-bin-decomposition}
\end{equation}
Then
$\sum_a\widehat\pi_a(b_a^2+r_a^2)=1$.
It is useful to use the normalized coordinates
\begin{equation}
 B_a=\sqrt{\widehat\pi_a}\,b_a,
 \qquad R_a=\sqrt{\widehat\pi_a}\,r_a,
 \qquad \sum_a(B_a^2+R_a^2)=1.
 \label{rcut:eq:comp-normalized-parameters}
\end{equation}
These coordinates place all bin parameters in a compact subset of a
unit sphere of fixed dimension.  They also accommodate vectors
concentrated on very few sites.  For example, if $p$ is one coordinate
vector in a bin of $m=n\widehat\pi_a$ sites, then
\[
 B_a=\frac1{\sqrt m},\qquad R_a=\sqrt{1-\frac1m},
\]
and all other coordinates vanish.  If instead $p$ is constant and
normalized on that bin, then $b_a=\widehat\pi_a^{-1/2}$, but
$B_a=1$ and $R_a=0$.

For any probability measure $\nu$ with finite second moment, denote
by $\mathcal V_{\beta,\nu}$ the expression in
Equation~\eqref{rcut:eq:variational} with every bracket and variance
taken under $\nu$, and with $D$ defined using that same measure.
Write $\mathcal A_{\beta,\nu}$ for its first two lines, so that its
last line is the remaining Gaussian comparison term.
For the empirical bin measure
$\widehat\nu_{\mathcal I}=\sum_a\widehat\pi_a\delta_{h_a}$,
let $\overline A=\diag(v_{a(i)})$, where $a(i)$ is the bin of $H_i$.
Equations~\eqref{rcut:eq:comp-reduced-form} and
\eqref{rcut:eq:comp-empirical-events} imply, with any fixed prescribed
accuracy and exponentially small failure,
\begin{equation}
 \sup_{\|p\|_2=1}
 \left|\widehat M(p)-
   \left[p^{\mathsf T}QWQ\overline A p
       +\mathcal A_{\beta,\widehat\nu_{\mathcal I}}(b,r)\right]
 \right|\ \text{is small}.
 \label{rcut:eq:comp-bin-form-error}
\end{equation}
{\emergencystretch=3em
The bin decomposition explains each part of this expression.
The diagonal term is
$\sum_a\widehat\pi_aD_{\widehat\nu_{\mathcal I}}(h_a)(b_a^2+r_a^2)$,
whereas terms bilinear in bin-constant vectors depend only on $b$.
The three rank terms in
Equation~\eqref{rcut:eq:comp-field-decomposition} give the three
$h$-terms of Equation~\eqref{rcut:eq:variational}, after replacing
$\overline H$ by $\beta^2$.\par}

Uniformity follows from the unit norm of $p$, without a pointwise
bound on $b$.  For instance, write
$\|f\|_{n,2}^2=n^{-1}\sum_i f_i^2$ and apply Cauchy--Schwarz:
\[
 \frac{|p^{\mathsf T}(H-H^{\mathcal I})|}{\sqrt n}
       \le\|H-H^{\mathcal I}\|_{n,2},
\]
and the corresponding error for $Hv$ is at most this error plus
$\max_i|v(H_i)-v_{a(i)}|\,\|H^{\mathcal I}\|_{n,2}$.
The other factors in the rank products are controlled by $\|p\|_2=1$
and the empirical second moment of the fields.  A uniform kernel
error contributes at most a constant times $\|p\|^2$ in the diagonal
term, and a constant times
$(n^{-1/2}\sum_i|p_i|)^2\le1$ in the cross term.  Finally, the
matrix term changes by at most
$\|W\|\max_i|v(H_i)-v_{a(i)}|$.  This proves the uniform error
estimate in Equation~\eqref{rcut:eq:comp-bin-form-error}.

We have now reduced the form to finitely many field values.  To apply
the hypothesis of Proposition~\ref{rcut:prop:planted-comparison}, we must
also replace empirical bin masses by their Gaussian probabilities.
For parameters $b_a,r_a$ normalized under $\widehat\pi$, define
functions on the original intervals by
\begin{equation}
 b^{\mathcal I}(h)=\frac{B_a}{\sqrt{\pi_a}},\qquad
 r^{\mathcal I}(h)=\frac{R_a}{\sqrt{\pi_a}}
       \quad(h\in I_a).
 \label{rcut:eq:comp-mass-rescaling}
\end{equation}
They satisfy $\langle(b^{\mathcal I})^2+(r^{\mathcal I})^2\rangle=1$.
By choosing the partition accurately enough and the proportion
tolerance small enough, we have
\begin{equation}
 \sup_{\sum_a(B_a^2+R_a^2)=1,\ R_a\ge0}
 \left|\mathcal V_{\beta,\widehat\nu_{\mathcal I}}(b,r)
       -\mathcal V_\beta(b^{\mathcal I},r^{\mathcal I})\right|
       \ \text{arbitrarily small}.
 \label{rcut:eq:comp-functional-approximation}
\end{equation}
The only nonpolynomial term requires care when its radical vanishes.
For bin probabilities $\pi$, set $t=(\sqrt{\pi_a})_a$ and
$c_{ad}=v_a+v_d$.  In normalized coordinates, write the $a$th
radical as a Euclidean norm:
\begin{equation}
 \mathsf N_a(t,B,R)=
 \left\|\left(\diag(c_{a1},\ldots,c_{ak})R,
       (I-tt^{\mathsf T})\diag(c_{a1},\ldots,c_{ak})B\right)\right\|_2
       \le2.
 \label{rcut:eq:comp-normalized-radical}
\end{equation}
If $t$ changes to another unit probability square-root vector $s$,
then $|\mathsf N_a(t,B,R)-\mathsf N_a(s,B,R)|\le4\|t-s\|_2$, since
$\|tt^{\mathsf T}-ss^{\mathsf T}\|\le2\|t-s\|_2$.
Consequently the last line of the functional, which is
$\beta\sum_at_aR_a\mathsf N_a(t,B,R)$, changes by at most
$6\beta\|t-s\|_2$.
The remaining terms are polynomials in $B,R$, the probability square
roots and the fixed bin coefficients.  Their coefficients are bounded,
so their continuity is uniform on the normalized sphere.  Since the
finitely many $\pi_a$ are positive, we can choose a sufficiently small
tolerance for the empirical proportions even when some fixed bins
have very small probability.

It remains to replace each representative by the actual field inside
its interval.  Again use a norm representation for the last line:
at outer field $h$, its radical is
\[
 \left\|\left((v(h)+v(\cdot))r(\cdot),
        (I-\E)((v(h)+v(\cdot))b(\cdot))\right)\right\|_{L^2\oplus L^2}.
\]
Changing $v$ uniformly by at most $a$ changes this vector by norm
at most $2a\sqrt{\langle b^2+r^2\rangle}=2a$.
This estimate is valid also when either radical is zero.
The diagonal and kernel terms use uniform errors, and the coherent
terms use, for example,
\[
 |\langle b(h-h^{\mathcal I})\rangle|
       \le\|b\|_2\|h-h^{\mathcal I}\|_2,
 \quad
 \|hv-h^{\mathcal I}v^{\mathcal I}\|_2
       \le\|h-h^{\mathcal I}\|_2
             +\|h^{\mathcal I}\|_2\|v-v^{\mathcal I}\|_\infty.
\]
The norm estimate and these $L^2$ estimates complete
Equation~\eqref{rcut:eq:comp-functional-approximation}.  In particular,
the comparison holds on the whole normalized class, including functions
whose mass is concentrated on rare bins.

\subsection{The conditional Gaussian comparison}

We now bound the remaining Gaussian matrix term.  We use the
centered expected-supremum comparison recalled by
Vitale~\cite[Equation~(3)]{Vitale2000}; the proof below verifies the
increment inequality for the present constraints.

Condition on $H$ and on the chosen bin partition.  Independence in
Equation~\eqref{rcut:eq:comp-field-decomposition} leaves the law of
$QWQ$ unchanged.  For this conditional calculation, represent it with
a fresh GOE independent of $H$ and continue to call that matrix $W$.
Fix also the empirical proportions and $b_a,r_a$.  We optimize only
over the residual vectors $u_a$, subject to the zero-sum and norm
conditions in Equation~\eqref{rcut:eq:comp-bin-decomposition}.

\begin{lemma}[Comparison within fixed bins]
 \label{rcut:lem:comp-bin-gaussian}
 With this conditioning and these fixed parameters,
 \begin{align}
 &\E\left[\sup_u
       p^{\mathsf T}QWQ\overline A p\,\middle|\,H\right]
 \notag\\[-2pt]
 &\qquad\le\beta\sum_a\widehat\pi_a r_a
    \sqrt{\sum_d\widehat\pi_d(v_a+v_d)^2r_d^2
           +\Var_{d\sim\widehat\pi}((v_a+v_d)b_d)}.
 \label{rcut:eq:comp-bin-gaussian-bound}
 \end{align}
 The full supremum on the left has Gaussian Lipschitz constant at
 most $\sqrt2\beta/\sqrt n$, uniformly in the normalized parameters.
\end{lemma}

\begin{proof}
 Separate the bin-constant space $\mathcal C$ from the zero-sum space
 $\mathcal U_a$ in each bin.  Then $\mathcal C$ is orthogonal to
 $\bigoplus_a\mathcal U_a$, and $\overline A$ preserves each space.
 Since $Qu=u$ and $Q\overline A u=\overline A u$, the process expands as
 \begin{align*}
 p^{\mathsf T}QWQ\overline A p
  ={}&C_{b}+
       \sum_a u_a^{\mathsf T}WQ(v_aI+\overline A)p_c
       +u^{\mathsf T}W\overline A u,\\
 C_b={}&p_c^{\mathsf T}QWQ\overline A p_c.
 \end{align*}
 For these fixed bin parameters, $C_b$ is constant in the residual
 index and has conditional mean zero.  The linear residual term uses
 the $\mathcal U$--$\mathcal C$ block of the GOE, while the quadratic
 term uses the $\mathcal U$--$\mathcal U$ block.  Their independence
 follows from Equation~\eqref{rcut:eq:comp-goe-covariance}.

 For two indices $u,u'$, set
 \[
  s_a=\widehat\pi_a r_a^2,
  \qquad q_a=u_a^{\mathsf T}u'_a,
  \qquad V_a=\Var_{d\sim\widehat\pi}((v_a+v_d)b_d).
 \]
 The linear residual covariance is
 \begin{equation}
  \frac{\beta^2}{n}\sum_a q_a V_a.
  \label{rcut:eq:comp-linear-covariance}
 \end{equation}
 Here $Q(v_aI+\overline A)p_c$ lies in $\mathcal C$ and has squared
 norm $V_a$.  Orthogonality to the residual spaces makes the other
 covariance contraction vanish.  Applying the same covariance identity
 to the quadratic residual term gives
 \begin{equation}
  \frac{\beta^2}{2n}\sum_{a,d}(v_a+v_d)^2q_aq_d.
  \label{rcut:eq:comp-quadratic-covariance}
 \end{equation}
 This follows from
 \[
  \frac{\beta^2}{n}
   \bigl((u^{\mathsf T}u')
          ((\overline A u)^{\mathsf T}\overline A u')
       +(u^{\mathsf T}\overline A u')
          ((\overline A u)^{\mathsf T}u')\bigr)
 \]
 and symmetrization in the bin indices.

 Cauchy--Schwarz gives $-s_a\le q_a\le s_a$, and hence
 \begin{equation}
  s_as_d-q_aq_d
       \le s_a(s_d-q_d)+s_d(s_a-q_a),
  \label{rcut:eq:comp-overlap-inequality}
 \end{equation}
 because the right side minus the left side is
 $(s_a-q_a)(s_d-q_d)\ge0$.  The inequality therefore holds also for
 negative overlaps.  Combining it with
 Equations~\eqref{rcut:eq:comp-linear-covariance}
 and~\eqref{rcut:eq:comp-quadratic-covariance} bounds the increment
 variance of the residual process by
 \[
  \frac{2\beta^2}{n}\sum_a(s_a-q_a)
       \left(V_a+\sum_d(v_a+v_d)^2s_d\right).
 \]
 This is exactly the increment variance of the linear process
 \[
  \frac\beta{\sqrt n}\sum_a
       \sqrt{V_a+\sum_d(v_a+v_d)^2s_d}\ z_a^{\mathsf T}u_a,
 \]
 with independent standard Gaussian vectors $z_a$ in $\mathcal U_a$.
 Lemma~\ref{rcut:lem:gaussian-facts} therefore bounds the expected
 supremum of the residual process by that of the linear process.
 We can maximize the latter separately in each residual space: each
 maximum is the corresponding Gaussian norm multiplied by
 $\sqrt{s_a}$.  Using
 $\E\|z_a\|\le\sqrt{\dim\mathcal U_a}
                     \le\sqrt{n\widehat\pi_a}$,
 we obtain Equation~\eqref{rcut:eq:comp-bin-gaussian-bound}.
 Adding $C_b$ leaves this expectation unchanged because its conditional
 mean is zero and the bin parameters remain fixed.

 Concentration must, however, include $C_b$.  For a unit vector $p$,
 the full GOE functional has symmetric coefficient matrix
 \[
  \tfrac12\bigl((Qp)(Q\overline A p)^{\mathsf T}
                +(Q\overline A p)(Qp)^{\mathsf T}\bigr),
 \]
 whose Frobenius norm is at most one.  Consequently its coefficient
 norm in independent standard Gaussian coordinates is at most
 $\sqrt2\beta/\sqrt n$.  The supremum inherits this Lipschitz
 constant, proving the assertion for the full process, including $C_b$.
\end{proof}

The comparison so far fixes $b$ and $r$.  We must now make it
simultaneous over these parameters.  Although $C_b$ has mean zero for
each fixed $b$, optimizing $b$ need not preserve that property.
We therefore use concentration for the full supremum just proved,
and pass through a finite net of the normalized parameters.

Fix a mesh size $\eta>0$.  The feasible set in
Equation~\eqref{rcut:eq:comp-normalized-parameters} has an $\eta$-net
inside that set of size at most $(1+2/\eta)^{2k}$, with an inessential
enlargement of this bound if necessary.  Conditional on $H$, this
net is fixed independently of $W$.  At each net point, adding the
deterministic $\mathcal A$ to the full supremum and applying
Lemma~\ref{rcut:lem:gaussian-facts} gives an upper bound
\begin{equation}
 \mathcal V_{\beta,\widehat\nu_{\mathcal I}}(b,r)+a
 \label{rcut:eq:comp-net-bound}
\end{equation}
outside probability $2\exp(-c_\beta n a^2)$, for every fixed $a>0$.
A union over the fixed net makes this estimate simultaneous at all
its points.

To pass from the net to every parameter, use the quadratic form in
Equation~\eqref{rcut:eq:comp-bin-form-error}.  Whenever $\|W\|$ and the
empirical field second moment are bounded, its operator norm is at most
a fixed $K_{\rm form}$.  Indeed, $\|W\|$ controls the random term,
the coherent terms are bounded rank forms, and the remaining kernels
are bounded.  The estimates above give exponentially small probability
for the complement of this event.  For any residual directions, changing
$(B,R)$ to $(B',R')$ while keeping the same unit direction in each
residual space changes the vector $p$ by exactly
$\|(B,R)-(B',R')\|_2$.  Hence the form changes by at most
$2K_{\rm form}\|(B,R)-(B',R')\|_2$.
If a residual norm is zero, choose any direction in its residual
space for this comparison.  The event giving at least two sites per
bin ensures that every such space is nonzero for large $n$.  Thus
the maximum of the net bounds, enlarged by $2K_{\rm form}\eta$,
controls the supremum over all normalized parameters.

\begin{proof}[Proof of Proposition~\ref{rcut:prop:planted-comparison}]
 We assemble the estimates with five fixed error allowances.  Given
 $\delta$ from the hypothesis, put $a=\delta/10$.  First require an
 error at most $a$ in Equation~\eqref{rcut:eq:comp-reduced-error}.
 Next choose the field partition and tolerances to make each error
 in Equations~\eqref{rcut:eq:comp-bin-form-error} and
 \eqref{rcut:eq:comp-functional-approximation} at most $a$.
 Take concentration error $a$ in Equation~\eqref{rcut:eq:comp-net-bound}
 and choose $\eta$ so that $2K_{\rm form}\eta\le a$.
 Every choice precedes the limit in $n$.
 The rescaling in Equation~\eqref{rcut:eq:comp-mass-rescaling} makes
 the limiting functions admissible for the hypothesis.  Functional
 approximation and concentration therefore bound each net value by
 $1-\delta+2a$.  The two form errors and extension from the net add
 the remaining $3a$, giving, for every unit vector,
 \[
  M_J(p)\le1-\delta+5a=1-\frac\delta2.
 \]

 To verify that this simultaneous assertion retains an exponential
 rate, collect the failure probabilities after fixing all parameters.
 Their sum has the form
 \begin{align*}
  &C n^2e^{-c_1n}+2nN_we^{-c_2n}
      +Ck e^{-c_3n}+2N_{\rm par}e^{-c_4n}
      +2\,9^ne^{-c_5n^2},\\
  &\hspace{25mm}N_w<\infty,\qquad
      N_{\rm par}\le(1+2/\eta)^{2k}.
 \end{align*}
 Each $c_i$ is positive and depends only on $\beta$ and the fixed
 choices.  The initial terms account for the entry, diagonal, norm,
 empirical-tail and common-shift events; $N_w$ counts the deterministic
 row-parameter mesh.  Only the last term uses an exponentially large
 sphere net, and its rate quadratic in $n$ absorbs $9^n$.  The other
 prefactors are fixed or polynomial in $n$.  Hence the sum is at most
 $e^{-\kappa n}$ for some $\kappa>0$ and all sufficiently large $n$.
 Finally the event in Equation~\eqref{rcut:eq:comp-exponential-conclusion}
 does not involve the auxiliary diagonals, so its probability on the
 enlarged space equals that on the planted marginal.
\end{proof}

\subsection{From planted probability to Gibbs mass}

The planted estimate controls a disorder event for one configuration.
Sign conjugation turns it into a bound on the expected unnormalized
weight of all exceptional configurations.  To obtain Gibbs mass, we
also need the partition function to remain near its annealed value
on the exponential scale.

\begin{lemma}[Annealed transfer]
 \label{rcut:lem:annealed-transfer}
 Fix $0<\beta<1$.  Let $\mathcal A_n$ be a measurable event for
 symmetric zero-diagonal interaction matrices, and suppose that
 $\widehat\P_n(\mathcal A_n)\le e^{-\kappa n}$ for all large $n$,
 with a fixed $\kappa>0$.
 For $x\in\{-1,1\}^n$, write
 $J^x=\diag(x)J\diag(x)$ and define
 \[
  \mathcal B_J=\{x:J^x\in\mathcal A_n\}.
 \]
 There are events whose original disorder probabilities tend to one
 on which
 \begin{equation}
  \mu_J(\mathcal B_J)\le e^{-\kappa n/2}.
  \label{rcut:eq:comp-gibbs-small}
 \end{equation}
 In fact the complementary probabilities can be bounded by
 $e^{-c n}$ for a fixed $c>0$.
\end{lemma}

\begin{proof}
 Let $Z_J$ be the original partition function and
 \[
  Z_{\mathcal B,J}=\sum_x e^{x^{\mathsf T}Jx/2}\1_{\{J^x\in\mathcal A_n\}}.
 \]
 The original disorder law is invariant under sign conjugation.  Each
 summand therefore reduces to the all-positive configuration, and the
 planted density in Equation~\eqref{rcut:eq:comp-planted-density} gives
 \begin{equation}
  \E Z_{\mathcal B,J}
    =2^n\E\left[e^{\sum_{i<j}J_{ij}}\1_{\mathcal A_n}\right]
    =\E Z_J\,\widehat\P_n(\mathcal A_n)
    \le e^{-\kappa n}\E Z_J.
  \label{rcut:eq:comp-annealed-numerator}
 \end{equation}
 Here $\E Z_J=2^n\exp(\beta^2(n-1)/4)$.

 We next bound the denominator $Z_J$ from below on the exponential
 scale.  Gaussian integration, with an independent uniform sign vector
 $(\varepsilon_i)_{i=1}^n$, gives the exact second-moment identity
 \begin{equation}
  \frac{\E Z_J^2}{(\E Z_J)^2}
    =e^{-\beta^2/2}\E_{\varepsilon}
       \exp\left\{\frac{\beta^2}{2n}
                    \left(\sum_i\varepsilon_i\right)^2\right\}
    \le\frac{e^{-\beta^2/2}}{\sqrt{1-\beta^2}}=:C_\beta.
  \label{rcut:eq:comp-partition-second-moment}
 \end{equation}
 For the inequality, introduce a standard Gaussian $Z$ to linearize
 the square.  The sign average becomes
 $\cosh(\beta Z/\sqrt n)^n\le e^{\beta^2Z^2/2}$, with expectation
 $(1-\beta^2)^{-1/2}$.  Separate the event
 $\{Z_J\ge\E Z_J/2\}$ and apply Cauchy--Schwarz to obtain
 \[
  \P_n(Z_J\ge\E Z_J/2)\ge\frac1{4C_\beta}.
 \]

 We upgrade this fixed positive probability using concentration.
 Each derivative of $\log Z_J$ in an independent standard disorder
 Gaussian has magnitude at most $\beta/\sqrt n$, giving Lipschitz
 constant at most $\beta\sqrt{(n-1)/2}$.  The preceding lower event
 and Lemma~\ref{rcut:lem:gaussian-facts} force
 \[
  \E\log Z_J\ge\log\E Z_J-C\sqrt n.
 \]
 Otherwise $\log\E Z_J-\log2$ would exceed the mean by a sufficiently
 large multiple of $\sqrt n$, and upper-tail concentration would make
 the preceding event less likely than its fixed lower bound.
 Lower-tail concentration now gives, for all large $n$,
 \begin{equation}
  \P_n\left(\log Z_J<\log\E Z_J-\frac{\kappa n}4\right)
      \le e^{-c n}.
  \label{rcut:eq:comp-partition-lower}
 \end{equation}
 Constants here may depend on the fixed $\beta$ and $\kappa$.

 By Equation~\eqref{rcut:eq:comp-annealed-numerator} and Markov's
 inequality,
 \[
  \P_n\left(Z_{\mathcal B,J}>
          e^{-3\kappa n/4}\E Z_J\right)\le e^{-\kappa n/4}.
 \]
 On the complement of this event and the event in
 Equation~\eqref{rcut:eq:comp-partition-lower}, divide the numerator
 bound by the partition-function lower bound to obtain
 Equation~\eqref{rcut:eq:comp-gibbs-small}.  A smaller positive exponent
 absorbs the sum of the two failure probabilities.
\end{proof}
